% page 52 of the facsimile (image p-051 of the scan)
% block 165.1 x 242.0 mm ; left margin 8.0 mm
\pgc{-12.700mm}{0.000mm}{
\l{\cel{6}44}
\l{}
\l{}
\l{\sou{asimptoto}. (\sou{geom.})\cel{2}Lineo rekta, ligata konstanta-relate a kurvo}
\l{\cel{3}di qua lu proximigas su sen atingar lu. - DEFIRS.}
\l{}
\l{\sou{asinektika}. (\sou{matem.}) (\sou{aludante anguli e figuri geometriala}). Ne-}
\l{\cel{3}homologa. - DEFIS.}
\l{}
\l{\sou{asistar}.\cel{2}(\sou{trans.})\cel{3}Esar prezenta ye... e regardar. - DFIRS.}
\l{}
\l{\sou{askarido}. (\sou{zool.}) Vermo hematoida qua vivas en la intestini di}
\l{\cel{3}la vertebrozi, precipue di la homo. - L.\cel{1}\sou{ascaridae genus :}}
\l{\cel{3}\sou{ascaris}. - DEFIRS.}
\l{}
\l{\sou{asketo}. I. Persono qua, en la komenco di la Eklezio kristana,}
\l{\cel{3}retroiris de la socio a loko ube lu esis sola, por ibe}
\l{\cel{3}konsakrar su ad exerci di pieso ed ad auto-punisi. - II.}
\l{\cel{3}Persono qua impozas a su vivo-maniero austera. -\cel{2}DEFIRS.}
\l{}
\l{\sou{asklepiado}.\cel{2}Verso qua konsistas ye un spondeo, du koriambi ed}
\l{\cel{3}un iambo. - DEFIRS.}
\l{}
\l{\sou{askoltar}.\cel{2}(\sou{trans.}) Atencar to quo dicesas ; esforcar por audar.}
\l{}
\l{}
\l{\sou{asno}.\cel{2}(\sou{zool.}) Animalo di la genero \textquotedbl{}solipedi\textquotedbl{}, kun oreli longa,}
\l{\cel{3}spino salianta, krio desharmonioza (bramo), e quan onu uzas}
\l{\cel{3}precipue kom kargajo-vestio. - L.\cel{1}\sou{asinus}.- EFIS.}
\l{}
\l{\sou{asociar}.\cel{2}(\sou{trans.}) Igar partoprenar to quon agas ulu. - Unionar}
\l{\cel{3}ad ulu kom partoprenonto di ta quon lu agas. - DEFIRS.}
\l{}
\l{\sou{asomar}.\cel{2}(\sou{trans.})Ocidar per frapo de ulo tre pezoza. - F.}
\l{}
\l{\sou{asonancar}.\cel{2}(\sou{netrans., kun}).\cel{2}I. Rivenar (\sou{aludante la sama son-}}
\l{\cel{3}\sou{uno)}. - II. Konsistar ye rimo nekompleta , quan karakterizas}
\l{\cel{3}nur la identeso di la vokalo acentoza inter du finali maskula}
\l{\cel{3}e du finali femina. (Ex. ek la Franca linguo :\cel{1}\sou{perte}\cel{1}e\cel{1}\sou{peste}).}
\l{\cel{3}- DEFIS.}
\l{}
\l{\sou{asortar}.\cel{2}(\sou{trans., kun}). Kunpozar kozi qui inter-harmonias.-DEFIR}
\l{}
\l{\sou{asparago}.\cel{2}(\sou{bot.}) Planto gardenala, di qua la yuna sprosi esas}
\l{\cel{3}manjebla. - L.\cel{1}\sou{asparagus}. - DEFIRS.}
\l{}
\l{\sou{asparagino}. (C8 H8 AZ2 O6). Substanco quan onu renkontras preci-}
\l{\cel{3}pue en la yuna sprosi di asparago e qua uzesas medicine}
\l{\cel{3}kom urinifiva. - DEFIRS.}
\l{}
\l{\sou{aspektar}.\cel{2}(netrans)Vidatesar kun karakterizivi qui igas judikar}
\l{\cel{3}lu kom... (ed esas ya tala). - DEFIRS.}
\l{}
\l{\sou{aspera}. Qua shokas desagreable la takto-senso. - EFIS.}
\l{}
\l{\sou{aspersar}.\cel{2}(\sou{trans., ye})\cel{2}Humidigar ye kelka guti. - dEFIRS.}
}
